% !TEX root = ../main.tex
\chapter{等式・書き換え・正規化}
\label{chap:ch12_equations_rewriting_normalization}

\begin{chapterabstract}
本章では、Lean 4 における等式証明を、仕様変更とリファクタリングの安全性を確認する道具として扱う。中心になる道具は \texttt{rw}、\texttt{simp}、\texttt{calc} である。\texttt{rw} は既知の等式に基づく置換、\texttt{simp} は自明な整理による正規化、\texttt{calc} は人間がレビューしやすい等式変形のログとして読む。本章のサンプルでは、価格計算ロジックを小さくモデル化し、関数分割、補助関数の導入、手数料計算の括弧変更が意味を変えないことを Lean に確認させる。
\end{chapterabstract}

\begin{learninggoals}
\item 等式証明を「変更前後で同じ結果になること」の確認として読める。
\item \texttt{rw}、\texttt{simp}、\texttt{calc} の役割を区別できる。
\item 価格計算のような小さな業務ロジックを、等式として整理できる。
\item 正規化を、レビューと自動証明の両方に効く設計方針として使える。
\end{learninggoals}

\section{この章を読む理由}

実務で頻繁に起きる変更の多くは、外から見た意味を変えないことが期待される。関数を分割する。補助関数を導入する。処理を共通化する。括弧の位置を変える。不要な初期値処理を消す。これらはすべて、理想的には「同じ結果を返す別の書き方」への変更である。

Lean では、この主張を等式として書く。

\begin{center}
\begin{tabular}{lll}
\toprule
実務上の言い方 & Lean での形 & 本章で使う道具 \\
\midrule
変更前後で結果が同じ & \texttt{old x = new x} & \texttt{rfl}, \texttt{rw}, \texttt{simp} \\
既知の仕様で置き換える & \texttt{h : a = b} を使う & \texttt{rw [h]} \\
自明な整理を自動化する & \texttt{x + 0} を \texttt{x} にする & \texttt{simp} \\
レビュー可能に段階を書く & 等式を一行ずつつなぐ & \texttt{calc} \\
\bottomrule
\end{tabular}
\end{center}

\begin{intuitionbox}
等式証明は、単に「数学っぽい証明」を書く作業ではない。リファクタリングの変更差分に対して、「この行の変更はこの仕様に基づく置換である」「この整理はゼロを足しても変わらないという規則である」と説明する作業である。
\end{intuitionbox}

\FigurePlaceholder{fig-ch12-rewriting-flow.png}{0.92}{等式証明を、変更前の式から正規形へ進む変形ログとして読む}{fig:ch12-rewriting-flow}

\section{仕様証明の多くは等式になる}

仕様にはさまざまな形がある。不変条件、事前条件、事後条件、順序関係、到達可能性、エラーケースの網羅などである。その中でも、リファクタリングやデータ変換で最初に現れるのは等式である。

たとえば、次のような主張はすべて等式として表せる。

\begin{itemize}
\item 料金計算を分割しても、合計金額は変わらない。
\item 何もしない正規化処理を二回通しても、一回通した結果と変わらない。
\item 手数料を一つずつ足しても、先にまとめてから足しても同じである。
\item 古い関数を新しい関数に置き換えても、すべての入力で同じ出力になる。
\end{itemize}

\begin{softwarebox}
この章では、本番コードそのものではなく、計算ロジックの核だけを Lean に移す。実務コードには小数、丸め、通貨単位、外部設定、キャンペーン条件などが含まれるかもしれない。しかし、まずは自然数で表せる小さなモデルを作り、「この書き換えはどの範囲で安全か」を明確にする。
\end{softwarebox}

\section{最初のモデル：価格計算を小さく切り出す}

最初に、割引後の小計に税を足すだけの小さな価格計算を考える。ここでは \texttt{Nat} を使う。\texttt{Nat} は自然数であり、負の値は扱わない。実務の金額型としては単純すぎるが、等式証明の形を学ぶには十分である。

\begin{leanbox}
\texttt{legacyPrice} は既存実装、\texttt{refactoredPrice} は書き換え後の実装だと読む。両者の定義を展開すると同じ式になるため、最初の証明は \texttt{rfl} で終わる。
\end{leanbox}

\begin{listing}[htbp]
\begin{minted}{lean4}
namespace Ch12

-- 割引後の小計。Nat の引き算は 0 で止まる。
def subtotal (base discount : Nat) : Nat :=
  base - discount

-- 税や手数料を加える単純な処理。
def addTax (amount tax : Nat) : Nat :=
  amount + tax

def legacyPrice (base discount tax : Nat) : Nat :=
  addTax (subtotal base discount) tax

def refactoredPrice (base discount tax : Nat) : Nat :=
  (base - discount) + tax

#eval legacyPrice 1000 100 80  -- => 980
#eval refactoredPrice 1000 100 80  -- => 980

theorem legacy_eq_refactored
    (base discount tax : Nat) :
    legacyPrice base discount tax =
      refactoredPrice base discount tax := by
  rfl
\end{minted}
\caption{価格計算の最小モデル}
\label{lst:ch12-price-model}
\end{listing}

\texttt{rfl} は反射律のタクティクである。直感的には、「両辺を定義に従って展開すると同じ形になる」と Lean が確認できる場合に使う。上の例では、\texttt{legacyPrice} と \texttt{refactoredPrice} は見た目は違うが、定義を展開すれば同じ式になる。

\begin{warningbox}
\texttt{rfl} で証明できるからといって、その仕様が実務上十分だとは限らない。このモデルでは、丸め、税率、負の金額、通貨単位、上限・下限を扱っていない。本章で証明しているのは、あくまで「この小さなモデルの範囲では変更前後が同じ」という主張である。
\end{warningbox}

\section{\texttt{rw}：仕様に基づく置換}

\texttt{rw} は rewrite の略である。既に持っている等式を使って、証明ゴールの中にある式を置き換える。

\begin{definitionbox}
等式 \texttt{h : a = b} があるとき、\texttt{rw [h]} はゴール中の \texttt{a} を \texttt{b} に置き換える。逆向きに使いたい場合は \texttt{rw [← h]} と書く。
\end{definitionbox}

この見方は、仕様レビューにそのまま対応する。補題 \texttt{addTax\_zero} は、「税額が 0 なら、税を足しても金額は変わらない」という小さな仕様である。その仕様を使って、\texttt{normalizedPrice} を \texttt{subtotal} に書き換える。

\begin{listing}[htbp]
\begin{minted}{lean4}
-- 0 を足す税計算は、元の金額と同じ。
theorem addTax_zero (n : Nat) :
    addTax n 0 = n := by
  unfold addTax
  rw [Nat.add_zero]

-- 既存コードには、不要な 0 税計算が残っているとする。
def normalizedPrice (base discount : Nat) : Nat :=
  addTax (subtotal base discount) 0

-- その不要な処理を仕様に基づいて取り除く。
theorem normalizedPrice_eq_subtotal
    (base discount : Nat) :
    normalizedPrice base discount =
      subtotal base discount := by
  unfold normalizedPrice
  rw [addTax_zero]
\end{minted}
\caption{\texttt{rw} による仕様ベースの置換}
\label{lst:ch12-rw}
\end{listing}

ここで重要なのは、\texttt{rw} が「文字列置換」ではないことである。\texttt{rw} が使うのは、Lean が検査済みの等式である。つまり、置換には根拠がある。コードレビューでいえば、「この変更は \texttt{addTax\_zero} という仕様に基づく」と言える状態である。

\subsection{逆向きの書き換え}

ときには、より複雑な形へ戻したいこともある。たとえば、既存のインターフェースに合わせるため、一時的に \texttt{n} を \texttt{addTax n 0} として見たい場合がある。このときは、等式を逆向きに使う。

\begin{listing}[htbp]
\begin{minted}{lean4}
theorem subtotal_eq_normalizedPrice
    (base discount : Nat) :
    subtotal base discount =
      normalizedPrice base discount := by
  rw [normalizedPrice_eq_subtotal]
\end{minted}
\caption{等式を逆向きに使う}
\label{lst:ch12-rw-backward}
\end{listing}

この例では、ゴールの右辺 \texttt{normalizedPrice base discount} が \texttt{subtotal base discount} に書き換えられ、最後は同じ式同士の等式になる。\texttt{rw} は左辺だけでなく、ゴール全体の中から一致する部分を探して書き換える。

\section{\texttt{simp}：自明な整理を正規化として使う}

\texttt{rw} は、どの等式で書き換えるかを人間が明示する。一方、\texttt{simp} は Lean が知っている単純化規則をまとめて使う。\texttt{x + 0 = x}、\texttt{0 + x = x}、\texttt{if true then a else b = a} のような明らかな整理は、毎回手で書くと読みづらくなる。

\begin{intuitionbox}
\texttt{simp} は「式を読みやすい標準形へ寄せる道具」と考えるとよい。つまり、正規化のための自動整理である。実務のリファクタリングでも、不要な wrapper、空の追加、恒等処理を消して、比較しやすい形へ寄せることがよくある。
\end{intuitionbox}

\begin{listing}[htbp]
\begin{minted}{lean4}
-- この補題は simp が自動で使えるようにしておく。
@[simp] theorem addTax_zero_simp (n : Nat) :
    addTax n 0 = n := by
  unfold addTax
  simp

-- 何もしない正規化処理を、あえて関数として置く。
def normalizeCode (n : Nat) : Nat :=
  n + 0

def loadCode (n : Nat) : Nat :=
  normalizeCode (normalizeCode n)

@[simp] theorem normalizeCode_eq (n : Nat) :
    normalizeCode n = n := by
  unfold normalizeCode
  simp

theorem loadCode_normalized (n : Nat) :
    loadCode n = n := by
  simp [loadCode]
\end{minted}
\caption{\texttt{simp} で使う正規化規則を登録する}
\label{lst:ch12-simp}
\end{listing}

\texttt{@[simp]} は、この補題を \texttt{simp} の整理規則として使ってよい、という印である。上の例では、\texttt{normalizeCode} が何もしない処理であることを一度証明しておく。その後、二重に呼ばれていても \texttt{simp} がまとめて取り除く。

\begin{warningbox}
\texttt{simp} に何でも登録すればよいわけではない。向きの悪い規則や、式を大きくする規則を登録すると、証明が読みにくくなったり、単純化が期待通りに進まなくなったりする。本章では、「明らかに小さくなる」「正規形に近づく」規則だけを登録する。
\end{warningbox}

\section{\texttt{calc}：レビュー可能な証明ログ}

\texttt{simp} は短く強力だが、何が起きたかが見えにくいことがある。仕様レビューで重要な証明では、あえて \texttt{calc} で途中式を残す方がよい。

\begin{definitionbox}
\texttt{calc} は、等式や不等式を段階的につなぐための構文である。各行に「なぜその変形が正しいか」を書けるため、人間が読む証明ログとして使いやすい。
\end{definitionbox}

次の例では、二つの手数料を順に足す既存実装と、先に手数料をまとめてから足す新実装が同じであることを示す。根拠は足し算の結合律である。

\begin{listing}[htbp]
\begin{minted}{lean4}
def addFee (amount fee : Nat) : Nat :=
  amount + fee

def legacyWithTwoFees
    (base service shipping : Nat) : Nat :=
  addFee (addFee base service) shipping

def refactoredWithTwoFees
    (base service shipping : Nat) : Nat :=
  addFee base (service + shipping)

theorem twoFees_eq_calc
    (base service shipping : Nat) :
    legacyWithTwoFees base service shipping =
      refactoredWithTwoFees base service shipping := by
  calc
    legacyWithTwoFees base service shipping
        = (base + service) + shipping := by rfl
    _ = base + (service + shipping) := by
          rw [Nat.add_assoc]
    _ = refactoredWithTwoFees base service shipping := by rfl
\end{minted}
\caption{\texttt{calc} による段階的な等式証明}
\label{lst:ch12-calc}
\end{listing}

この証明は、同じ定義展開と結合律を \texttt{simp} に渡して、もっと短く書けるかもしれない。しかし、レビュー目的なら、\texttt{calc} の方が「どこで結合律を使ったか」が明確である。

\begin{softwarebox}
\texttt{calc} は、設計文書に残す証明として相性がよい。レビュー担当者は Lean の全タクティクを知らなくても、「既存実装を展開する」「結合律で括弧を変える」「新実装の定義に戻す」という流れを読める。
\end{softwarebox}

\section{正規化：比較しやすい形へ寄せる}

正規化とは、同じ意味を持つ複数の書き方を、比較しやすい代表形へ寄せることである。Lean の \texttt{simp} は、まさにこの作業を助ける。実務でも、比較しやすい形にしてから差分を見ると、レビューが楽になる。

\begin{definitionbox}
本章でいう正規形とは、次のような形である。
\begin{itemize}
\item 不要な \texttt{+ 0} や恒等処理が消えている。
\item 補助関数が必要な範囲だけ展開されている。
\item 括弧の付き方が一貫している。
\item 変更前と変更後を同じ構造で比較できる。
\end{itemize}
\end{definitionbox}

\FigurePlaceholder{fig-ch12-price-normalization.png}{0.90}{価格計算ロジックを同じ正規形へ寄せて比較する}{fig:ch12-price-normalization}

次の例では、既存コードが小計、二つの手数料、税額 0 の追加を段階的に行っている。新コードでは、税額 0 を消し、二つの手数料をまとめている。\texttt{simp} と結合律を使うと、両者は同じ正規形へ整理できる。

\begin{listing}[htbp]
\begin{minted}{lean4}
def chargeLegacy
    (base discount service shipping : Nat) : Nat :=
  addTax
    (addFee (addFee (subtotal base discount) service) shipping)
    0

def chargeNormalized
    (base discount service shipping : Nat) : Nat :=
  (base - discount) + (service + shipping)

theorem chargeLegacy_eq_chargeNormalized
    (base discount service shipping : Nat) :
    chargeLegacy base discount service shipping =
      chargeNormalized base discount service shipping := by
  unfold chargeLegacy chargeNormalized subtotal addFee addTax
  simp [Nat.add_assoc]
\end{minted}
\caption{価格計算ロジックの正規化}
\label{lst:ch12-price-normalization}
\end{listing}

この証明の読み方は次の通りである。

\begin{enumerate}
\item まず \texttt{unfold} で、業務名の付いた補助関数を計算式へ展開する。
\item 次に \texttt{simp} で、\texttt{+ 0} のような自明な部分を消す。
\item 最後に \texttt{Nat.add\_assoc} で、足し算の括弧を新実装の形へそろえる。
\end{enumerate}

ここでの \texttt{unfold} は、抽象化を一時的にほどく操作である。普段は \texttt{subtotal} や \texttt{addFee} の名前を使う方が読みやすい。しかし、変更前後が本当に同じか比較するときには、必要な範囲だけ定義を展開する。

\section{\texorpdfstring{\texttt{rw}、\texttt{simp}、\texttt{calc} の使い分け}{rw、simp、calc の使い分け}}

三つの道具は重なり合うが、目的は少し違う。

\begin{center}
\begin{tabular}{p{0.18\linewidth}p{0.34\linewidth}p{0.36\linewidth}}
\toprule
道具 & 向いている場面 & 実務上の読み方 \\
\midrule
\texttt{rw} & 特定の補題を根拠に一箇所を書き換える & この変更はこの仕様に基づく置換である \\
\texttt{simp} & 自明な整理をまとめて行う & 不要なノイズを消して正規形へ寄せる \\
\texttt{calc} & 途中式を残して説明したい & レビュー可能な変更ログを書く \\
\bottomrule
\end{tabular}
\end{center}

\begin{leanbox}
証明が短ければよいとは限らない。Lean に慣れた人だけが読む補題なら \texttt{simp} で十分なことがある。一方、仕様レビューに載せる証明では、\texttt{calc} で変形理由を明示した方が有用である。
\end{leanbox}

\section{よくある誤解}

\subsection{誤解1：\texttt{simp} は何でも証明してくれる}

\texttt{simp} は強力だが、万能ではない。基本的には、登録された単純化規則や定義展開によって式を整理する道具である。複雑な業務仕様を勝手に理解するわけではない。\texttt{simp} に使わせたい業務上の事実があるなら、それを補題として先に証明する必要がある。

\subsection{誤解2：\texttt{rw} はプログラムの文字列置換である}

\texttt{rw} はテキストエディタの置換ではない。Lean が型検査した等式だけを使う。置換してよい理由が証明オブジェクトとして存在するため、単なる検索置換よりずっと厳密である。

\subsection{誤解3：短い証明ほど良い証明である}

短い証明は保守しやすいことが多い。しかし、仕様説明としては途中式が必要な場合もある。特に設計レビューでは、\texttt{calc} による冗長な証明が、将来の読者にとって最も読みやすいことがある。

\section{実務へ持ち込むときの手順}

本章の考え方を実務のリファクタリングに使う場合、最初から大きなコードベース全体を Lean に移す必要はない。次の順で進めると扱いやすい。

\begin{enumerate}
\item 変更で守りたい性質を一文で書く。例：「手数料を先にまとめても合計は変わらない」。
\item その性質に関係する入力と出力だけを型として切り出す。
\item 変更前と変更後を、小さな Lean 関数として定義する。
\item 明らかな仕様を補題にする。例：\texttt{addTax n 0 = n}。
\item \texttt{rw}、\texttt{simp}、\texttt{calc} で等式を証明する。
\item 証明した範囲と、実装側に残る前提を文書に書く。
\end{enumerate}

\begin{warningbox}
Lean の証明モデルと本番実装は同じではない。特に、丸め、小数、タイムゾーン、外部設定、DB の NULL、例外処理は、単純な \texttt{Nat} モデルでは表せないことが多い。証明済みと言えるのは、Lean に書いたモデルと前提の範囲に限られる。
\end{warningbox}

\section{本章のまとめ}

リファクタリングや計算式の整理は、多くの場合「変更前後が等しい」という主張として表せる。

\texttt{rw} は、検査済みの等式を使う根拠付きの置換である。

\texttt{simp} は、不要なノイズを消して式を正規形へ寄せる道具である。

\texttt{calc} は、途中式と変形理由を残すレビュー可能な証明ログである。

正規化を意識すると、変更前後のロジックを比較しやすくなる。

Lean で証明した範囲と、実務コード側に残る前提は必ず分けて書く。

次章では、\texttt{Bool}、\texttt{Option}、\texttt{Sum} のような分岐を含む型を扱う。等式だけで済む例では \texttt{rw}、\texttt{simp}、\texttt{calc} が中心だった。分岐がある仕様では、すべてのケースを漏れなく扱うために \texttt{cases} が必要になる。
